\section{Return to the original fibration and consequences}
\label{application:section}

The previous sections supply the two positivity statements and the
exact section comparison. We now apply them to the original fibration,
keeping its base fixed while changing the intermediate Iitaka space.

\begin{proof}[Proof of \Cref{main:orbifold}]
If either summand on the right is \(-\infty\), the asserted inequality
is automatic.  Suppose henceforth that both are nonnegative, and put
\[
 k=\kappa(F,K_F+\Delta_F),\qquad b=\kappa(f\mid\Delta).
\]

\smallskip\noindent
\emph{Fixing the base.}
Use Section~\ref{setup:section} and \Cref{setup:flattening} to resolve the diagram,
take K\"ahler models, and flatten before resolving the source. Use
strict transforms plus reduced exceptional divisors at every source
modification. Resolve the discriminant and boundary-stratum images
on the base, then fix this smooth K\"ahler base \(Y\). The source
and very general fiber log section rings are unchanged, and every
source prime mapped into codimension at least two in \(Y\) has
coefficient one. Put
\[
 C=B(f,\Delta)
\]
on this model. It has SNC support, and the definition of the invariant
gives
\begin{equation}
 \kappa(Y,K_Y+C)\geq b.
 \label{application:invariant-bound}
\end{equation}
\Cref{setup:simultaneous} preserves \(C=B(f,\Delta)\) and the
coefficient-one condition through all further changes over \(Y\).

\smallskip\noindent
\emph{The relative Iitaka map and its Hodge line.}
Apply \Cref{application:relative-iitaka}, obtaining
\[
 X\xrightarrow{g}W\xrightarrow{h}Y,
 \qquad \dim W_y=k,
 \qquad \kappa(G,K_G+\Delta_G)=0.
\]
Make \(W\) K\"ahler, resolve the diagram, and fix the current smooth
source pair \((X_0,\Delta_0)\) as reference. Flatten over \(W\)
and resolve the generator and boundary data as required by
\Cref{rankone:comparison}. Thus every prime mapped by \(g\) into
codimension at least two is exceptional over \(X_0\), and on a dense
smooth log open the relative
systems have rank one in all sufficiently divisible degrees.

Apply \Cref{rankone:comparison,period:adjoint}. Their further
modifications remain over \(Y\); \Cref{setup:simultaneous} preserves
the exceptional reference, while the variation and its parabolic line
pull back from the common open. The relative Iitaka property persists
by \Cref{application:relative-iitaka}. On the resulting common model
we obtain a rational line \(M\), a rational SNC boundary \(T\) in
\([0,1]\), and a surjective map to a smooth projective variety:
\begin{equation}
 p:W\longrightarrow S,
 \qquad M=p^*P,
 \qquad P\text{ nef},
 \qquad K_S+a_0P\text{ big for some }a_0>0,
 \label{application:period-data}
\end{equation}
such that
\begin{equation}
 B(g,\Delta)\leq T\leq1.
 \label{application:rankone-boundary}
\end{equation}
Compute \(T\) by the rank-one order formula on this final model.
Parabolic pullback and \Cref{rankone:comparison} identify the divisible
systems of \(K_W+T+M\) with those on \((X_0,\Delta_0)\), also
relatively over \(Y\), preserving all ratios.

\smallskip\noindent
\emph{Relative bigness and the base boundary.}
Set \(D_W=K_W+T\).  The systems defining the relative Iitaka image
have image dimension \(k\) on very general \(f\)-fibers. Push these systems
into \(D_W+M\) using \Cref{rankone:comparison}.  Their ratios are
unchanged, so the resulting systems have image dimension \(k\) on
very general \(h\)-fibers, whose dimension is \(k\).  Consequently
\begin{equation}
 (D_W+M)|_{W_y}\text{ is big for very general }y.
 \label{application:relative-big}
\end{equation}
Here the systems are obtained by generic base change over an open patch
of \(Y\) and extend on the smooth \(W_y\) by the relative section
comparison. Adjunction identifies the restricted adjoint bundles.

By \Cref{application:composition},
\begin{equation}
 B(h,T)\geq C.
 \label{application:composed-boundary}
\end{equation}
The same lemma gives coefficient one in \(T\) for every prime of
\(W\) mapped into codimension at least two in \(Y\). Thus both
boundary hypotheses of \Cref{generaltype:addition} hold.

The hypotheses of \Cref{addition:principle} have now all been
verified. Its conclusion is
\begin{equation}\label{application:final-intermediate-bound}
 \kappa(W,K_W+T+M)\geq k+\kappa(Y,K_Y+C).
\end{equation}
The proof of that proposition uses a single nonzero section to cancel
the ample perturbation and preserves the ratios defining this bound.

Finally, \Cref{rankone:comparison} lifts these systems by the relative
generators, preserving their ratios and image dimensions. Untested
\(g\)-contracted primes are exceptional over \((X_0,\Delta_0)\),
so the forms descend and extend there. By
\Cref{setup:log-modification}, they then descend to the original source.
Combining \eqref{application:invariant-bound} and
\eqref{application:final-intermediate-bound} proves
\[
 \kappa(X,K_X+\Delta)
     \geq \kappa(F,K_F+\Delta_F)+\kappa(f\mid\Delta).
\]

Bigness on a point means Iitaka dimension zero, and a zero-dimensional
connected fiber has trivial relative line. If \(Y\) is a point the
inequality is an identity; if the relative Iitaka fiber is a point its
Hodge line is trivial. Together with the point-period-base case above
and the initial \(-\infty\) convention, these observations cover all
dimensions and values of the two summands.
\end{proof}

\subsection{The logarithmic and ordinary inequalities}
\label{spec:section}

We first compare the invariant orbifold base with a prescribed
logarithmic base.  This makes the two classical specializations of
Theorem~\ref{main:orbifold} immediate, while retaining control of the
base under modification.

\begin{lemma}[Comparison with a logarithmic base]
\label{spec:base-comparison}
Let $f:(X,\Delta)\to Y$ be a fibration of smooth compact complex
manifolds, where $\Delta$ has rational coefficients in $[0,1]$ and
simple normal crossing support.  Let $D_Y$ be a reduced simple normal
crossing divisor on $Y$.  Suppose that every prime component of
$f^*D_Y$ has coefficient one in $\Delta$.  Then
\[
  \kappa(f\mid\Delta)\geq\kappa(Y,K_Y+D_Y).
\]
\end{lemma}

\begin{proof}
Choose a neat model as in Lemma~\ref{setup:flattening}, with
$p:X'\to X$, $q:Y'\to Y$, $f':X'\to Y'$, and
$\Delta'=p_*^{-1}\Delta+\operatorname{Exc}(p)_{\mathrm{red}}$.
We may arrange that
$E_Y=(q^{-1}\operatorname{Supp}D_Y)_{\mathrm{red}}$ has simple
normal crossings.  Every prime of $X'$ dominating a component of
$E_Y$ has coefficient one in $\Delta'$: an exceptional prime has
coefficient one by definition, and a nonexceptional prime is the strict
transform of a component of $f^*D_Y$.  Thus its weighted multiplicity
is infinite, and the inf-multiplicity rule gives
\[
 B(f',\Delta')\geq E_Y.
\]
Pullback by $q$ takes logarithmic top forms with poles along $D_Y$
to logarithmic top forms with poles along $E_Y$.  In local coordinates
this follows by pulling back each factor $dz_i/z_i$; the same assertion
holds for all tensor powers.  Consequently, pullback embeds every
plurilogarithmic system on $(Y,D_Y)$ into the corresponding system of
$K_{Y'}+B(f',\Delta')$, preserving all ratios of sections.
The neat-model identity~\eqref{setup:neat-formula} now proves the
claim.
\end{proof}

\begin{corollary}[Logarithmic subadditivity]
\label{spec:logarithmic}
Let $f:X\to Y$ be a surjective morphism with connected fibers between
smooth connected projective complex varieties.  Let $D_X,D_Y$ be
reduced simple normal crossing divisors, with the zero divisor allowed,
such that
\[
 \operatorname{Supp}(f^*D_Y)\subseteq\operatorname{Supp}(D_X).
\]
For a very general fiber $F$, put $D_F=D_X|_F$.  Then
\begin{equation}\label{spec:log-inequality}
 \kappa(X,K_X+D_X)
 \geq\kappa(F,K_F+D_F)+\kappa(Y,K_Y+D_Y).
\end{equation}
The same statement holds for compact manifolds in Fujiki class
$\mathcal C$ and holomorphic fibrations between them.
\end{corollary}

\begin{proof}
Theorem~\ref{main:orbifold} applies to $(X,D_X)$, and
Lemma~\ref{spec:base-comparison} bounds its base term from below.
If either term on the right of~\eqref{spec:log-inequality} is
$-\infty$, the assertion follows from our convention for that value.
\end{proof}

In the projective setting, this is the usual compactification form of
the logarithmic Iitaka conjecture.  Indeed, for a smooth
quasi-projective variety $U$ with a smooth projective compactification
$X$ and reduced simple normal crossing boundary $D_X=X\setminus U$,
its logarithmic Kodaira dimension is
$\overline\kappa(U)=\kappa(X,K_X+D_X)$
\cite{IitakaLog}.  A dominant morphism $U\to V$ between smooth
connected quasi-projective varieties, with geometrically connected
general fiber, admits compatible such compactifications:
compactify the graph and resolve the graph and boundaries.  The
Stein factor of the resulting morphism $X\to Y$ is finite
birational over the normal variety $Y$, hence equals $Y$.
Thus $f$ has connected fibers, and
$\operatorname{Supp}(f^*D_Y)\subseteq\operatorname{Supp}(D_X)$.
For very general $v\in V$, $(X_v,D_X|_{X_v})$ compactifies $U_v$.
Thus Corollary~\ref{spec:logarithmic} gives
$\overline\kappa(U)\geq\overline\kappa(U_v)+\overline\kappa(V)$.

Ordinary subadditivity also enters the Albanese reduction of
\cite{OpenAILogAbundance2026}. That work combines
it with separate nonvanishing and good-model arguments to prove log
abundance in characteristic zero.

\begin{corollary}[Ordinary subadditivity in characteristic zero]
\label{spec:ordinary}
Let $k$ be an algebraically closed field of characteristic zero, and
let $f:X\to Y$ be a surjective morphism with connected fibers between
smooth connected projective $k$-varieties.  If $F$ is its geometric
generic fiber, then
\begin{equation}\label{spec:ordinary-inequality}
 \kappa(X)\geq\kappa(F)+\kappa(Y).
\end{equation}
More generally, the logarithmic inequality of
Corollary~\ref{spec:logarithmic}, with the same boundary hypotheses,
holds over $k$, with the geometric generic fiber in place of a very
general complex fiber.
\end{corollary}

\begin{proof}
Over $\mathbb C$, the ordinary assertion follows by taking
$D_X=D_Y=0$ in Corollary~\ref{spec:logarithmic}.  In a projective
family over $\mathbb C$, the Iitaka dimension of the geometric
generic fiber equals that of a very general fiber.  To see this,
apply cohomology and base change in each positive divisible
pluricanonical degree.  After shrinking the base separately in each
degree, the corresponding relative rational map also has constant
fiberwise image dimension.  Excluding the resulting countable
union of proper closed subsets identifies all these image
dimensions with those of the generic fiber.

For the field extension, let $K\subseteq K'$ be fields and let $V$
be a smooth projective geometrically integral $K$-variety with a
rational divisor $D$.  For every $m$ for which $mD$ is integral,
flat base change gives
\[
 H^0(V,\mathcal O_V(mD))\otimes_K K'
   \simeq H^0(V_{K'},\mathcal O_{V_{K'}}(mD_{K'})).
\]
The complete linear system on the right is the scalar extension of
the one on the left, so nonvanishing and image dimension are
unchanged.  Hence Iitaka dimension is invariant under extension
of the ground field.  Applying this observation to the generic
fiber over the function field of the base proves the corresponding
invariance for geometric generic fibers.

Descend $f$, its projective embeddings, and its boundary divisors
to a finitely generated subfield $K\subset k$ over $\mathbb Q$.
The descended varieties are smooth and geometrically integral;
the boundary conditions descend as well.  The equality
$f_*\mathcal O_X=\mathcal O_Y$ descends by proper flat base change
and faithful flatness, and therefore also holds after an embedding
$K\hookrightarrow\mathbb C$.  Apply the complex logarithmic
inequality to this base change.  The preceding invariance identifies
its three Iitaka dimensions with those over $k$, proving both
assertions.
\end{proof}

\subsection{Kodaira dimension along the core}
\label{spec:core}

A smooth compact orbifold pair $(X,\Delta)$ in class $\mathcal C$ is
\emph{special} if it has no meromorphic fibration onto a
positive-dimensional base of orbifold general type; here general
type means that the invariant base dimension equals the dimension
of the base.  Campana constructed its \emph{core}
\[
 c_{(X,\Delta)}:(X,\Delta)\dashrightarrow C(X,\Delta),
\]
whose very general orbifold fiber is special and whose orbifold base is of general
type, unless it is a point
\cite[Theorem~9.1]{CampanaOrbifolds}.  In particular, the core has
a point base exactly when $(X,\Delta)$ is special. Fiber properties
of the core and the tower below refer to Campana's stable orbifold
fibers: they are computed on the suitable equivalent holomorphic
representatives furnished by these constructions
\cite[Definition~2.49 and \S7.1]{CampanaOrbifolds}.
A neat representative $g:(Z,\Gamma)\to T$ in the sense of
Section~\ref{setup:section} is \emph{strictly neat} if its
inf-multiplicity boundary on the smooth base is SNC. It is
\emph{high} if $g$ is an orbifold morphism to this base: for every
base prime $D$ and source prime $E$ with
$a_E=\ord_E(g^*D)>0$, including primes contracted by $g$,
\[
              a_E m_\Gamma(E)\geq m_{B(g,\Gamma)}(D).
\]
These are the model conditions of
\cite[Definition~3.8 and Proposition~3.10]{CampanaOrbifolds}.
The high condition is not automatic: the inf-multiplicity base tests
divisors dominating base primes, whereas an orbifold morphism must
also control primes contracted into higher codimension
\cite[Remark~3.3(2)]{CampanaOrbifolds}.
For the core and its tower we choose compatible strictly neat and
high representatives supplied by Campana's constructions
\cite[Proposition~3.15 and Corollary~4.14]{CampanaOrbifolds}.
The section comparisons of Section~\ref{setup:section} do not
assert that arbitrary model changes preserve every stable-fiber
property.

\begin{corollary}[Kodaira dimension along the core]
\label{spec:core-dimension}
Let $(X,\Delta)$ be a smooth compact orbifold pair in Fujiki class
$\mathcal C$, with rational simple normal crossing boundary
coefficients in $[0,1]$.  Let $C=C(X,\Delta)$ be its core base and
$(G,\Delta_G)$ a very general orbifold fiber.  Then
\begin{equation}\label{spec:core-equality}
 \kappa(X,K_X+\Delta)
   =\kappa(G,K_G+\Delta_G)+\dim C.
\end{equation}
If $\kappa(X,K_X+\Delta)\geq0$, then
\[
 0\leq\dim C\leq\kappa(X,K_X+\Delta),
\]
and equality in the second inequality holds precisely when
$\kappa(G,K_G+\Delta_G)=0$.
In particular, $\kappa(X,K_X+\Delta)=0$ implies that
$(X,\Delta)$ is special.
\end{corollary}

\begin{proof}
Resolve the core map, using the boundary convention of
Lemma~\ref{setup:log-modification}.  The total and fiber Iitaka
dimensions are preserved.  Theorem~\ref{main:orbifold} gives
the lower bound in~\eqref{spec:core-equality}, since the invariant
base dimension is $\dim C$.

For completeness, the reverse bound is the elementary upper
addition inequality.  For each divisible $m$, let $\phi_m$ be
the meromorphic map defined by $|m(K_X+\Delta)|$, when this
system is nonempty.  The map $(c,\phi_m)$ has image dimension
at most $\dim C+\kappa(G,K_G+\Delta_G)$: on a very general
$c$-fiber its second component is given by a subspace of
$H^0(G,m(K_G+\Delta_G))$.  Since projection from this image
onto the image of $\phi_m$ is surjective, the same upper bound
holds for $\dim\phi_m(X)$.  Taking the maximum over $m$ proves
the desired inequality.  If the fiber Iitaka dimension is
$-\infty$, restriction to a very general fiber shows that no
global pluricanonical section can be nonzero, so both sides of
\eqref{spec:core-equality} are $-\infty$.

Finally, if the total Iitaka dimension is nonnegative, a nonzero
section restricts nontrivially to a very general core fiber.
Its Iitaka dimension is therefore nonnegative.  The dimension
bound and its equality case follow from~\eqref{spec:core-equality}.
When the total dimension is zero, the core base is a point,
which is equivalent to specialness.
\end{proof}

These statements recover, for the core, Campana's general-type-base addition
theorem and its consequences
\cite[Theorems~6.3, 6.7, and~6.8]{CampanaOrbifolds}.
The existence of the core and the implication from Kodaira
dimension zero to specialness are established results of his
theory.

\subsection{The core tower}
\label{spec:core-tower-section}

For a smooth orbifold pair $(Z,\Theta)$, the notation
$\kappa_+(Z,\Theta)=-\infty$ means that every meromorphic
fibration onto a positive-dimensional base has invariant
orbifold Kodaira dimension $-\infty$; the condition is vacuous
for a point.  Equivalently, one may test all dominant meromorphic
maps: in Stein factorization, the finite orbifold canonical-divisor
inequality shows that passing to the connected-fiber factor cannot
decrease the invariant base dimension.  Indeed, over a divisor
with multiplicity $m_D$, a prime with ramification index $e$ and
orbifold multiplicity $m_E$ satisfies $m_D\leq e m_E$; the
canonical-divisor difference has coefficient
$e/m_D-1/m_E\geq0$.
The \emph{$\kappa$-rational quotient} is a meromorphic fibration
$r:(Z,\Theta)\dashrightarrow R$ with $\kappa(r\mid\Theta)\geq0$
whose very general stable orbifold fiber has $\kappa_+=-\infty$
\cite[Definition~5.23 and Corollary~6.14]{CampanaOrbifolds}.
Its \emph{stable orbifold base} is the inf-multiplicity base on the
compatible representatives just described
\cite[Proposition~3.15, Corollary~4.14, and \S7.1]{CampanaOrbifolds}.

\begin{corollary}[Campana's decomposition of the core]
\label{spec:core-tower}
Let $(X,\Delta)$ be as in Corollary~\ref{spec:core-dimension},
and set $n=\dim X$.  Starting from $(X,\Delta)$, successively
take the $\kappa$-rational quotient $r$ and the
Moishezon--Iitaka fibration $M$ of its stable orbifold base.
At each iteration choose such compatible representatives and equip
the new base with the stable inf-multiplicity orbifold structure
of the composite.
The composite after at most $n$
iterations is the core, up to bimeromorphic equivalence of
fibrations:
\[
 c_{(X,\Delta)}=(M\circ r)^n.
\]
The positive-dimensional very general stable orbifold fibers of
the $r$-steps have $\kappa_+=-\infty$, and those of the
$M$-steps have Kodaira dimension zero.  In particular, for
a special pair this tower ends at a point.
\end{corollary}

\begin{proof}
Theorem~\ref{main:orbifold} establishes the hypothesis of
Campana's Corollary~6.14, which supplies the
$\kappa$-rational quotient.  Its scope is exactly that used
here: smooth compact class-$\mathcal C$ pairs with rational
coefficients in $[0,1]$
\cite[Conjecture~6.1 and Remark~6.2]{CampanaOrbifolds}.
Its stable base has nonnegative Kodaira dimension, so $M$
is defined and has very general orbifold fiber of Kodaira
dimension zero.  The new smooth base remains in class
$\mathcal C$, with rational boundary coefficients in $[0,1]$.
Thus the construction can be repeated.

Campana's Lemma~10.1 states that each composite $M\circ r$
has special general orbifold fiber and decreases dimension
unless its source pair is of general type.  His
Theorems~10.2--10.3 identify the resulting stabilized
composite with the core.  Their sole conjectural input is
the existence of the $\kappa$-rational quotient, now supplied
above.  Strict dimension decrease allows at most $n$
nontrivial iterations; identity steps give the displayed
formula.  A point is fixed throughout.
These are precisely the decomposition and fiber assertions of
\cite[Lemma~10.1, Theorems~10.2--10.3, and Corollary~10.4]
{CampanaOrbifolds}.
\end{proof}
